\section{Conditions des réalisations}
\label{sec:ii-condiciones-realizaciones}

Les réalisations d’action, de mélange, d’aire et de gravitation conservent
les domaines et les normalisations des résultats qui les composent.

\subsection{Coordonnée d’alpha et repères de mélange}

La clôture dodécaphasique et la représentation analytique d’alpha expriment
la même coordonnée à toute profondeur
(théorème~\ref{thm:alpha-dos-publicaciones-completas}).
La récurrence analytique est fixée à partir de l’état enrichi et de la classe
géométrique déclarée ; elle ne sélectionne pas indépendamment cette coordonnée.
Dans la réalisation angulaire, les repères spectraux complets interviennent
dans \eqref{eq:ii-pmns-frames}. L’unitarité d’un facteur polaire ne remplace pas
son identification à ce changement de repères. La compatibilité CP conserve
la condition \eqref{eq:ii-cp-square}, et une forme réelle de Majorana restreint
les phases admissibles.

\subsection{Bases thermiques et normalisation gravitationnelle}

La section~\ref{sec:ii-boltzmann} définit le transducteur thermique et sa
transformation sous changement de bases. Le produit
\(k_B\Theta_{\rm clk}\log3\) est fixé par l’énergie du cycle ;
la coordonnée individuelle de \(k_B\) exige de spécifier les bases positives
des droites d’énergie et de température. Cette covariance n’altère
ni l’action ni la période précédentes.

La section~\ref{sec:gravity} fixe la coordinatisation gravitationnelle au moyen du rayon
réduit \(Gm/c^2\) et de sa condition de coïncidence avec \(R_{108}\).
Dans la composition actuelle, la longueur précède cette expression circulaire :
R4 et R5 déterminent \(L\) et le rayon positif du caractère, et l’horloge
\(t_0=\pi_{\HMT}L/(54c_{\rm int})\) donne \(R_{108}=L\).
R8 identifie ce rayon à la lecture gravitationnelle réduite.
Les sections d’action réciproques conservent le produit \(\hbar G\)
et, par conséquent, l’échelle aréale correspondante. L’équation de Cartan--Holst
se formule pour une cotétrade non dégénérée, une connexion et le courant
de spin déclarés, avec les conditions aux limites de sa variation.

\subsection{Provenance de la composition longitudinale et radiale}

La conservation de l’archive, le transport de Hadamard, l’action de
phase et la réciprocité des sections proviennent des développements
antérieurs du corpus. La révision des 25--26 septembre 2026
réunit le registre marqué d’incidences et explicite les compositions
longitudinale R4 et métrique R5. Le théorème de
\ref{g26:sec:rigidez-radial-es} démontre leur conséquence conjointe avec
R1--R3 et R6--R8 : unicité de l’inscription, du lecteur radial et du
coefficient commun dans la réalisation déclarée \cite{hmtg26radial}.
Cette provenance distingue les antécédents conservés des
définitions et preuves incorporées dans la révision.

Les transports avec métrique globale, y compris leurs blocs croisés,
conservent l’action et le résidu récupérable. La restitution du
contre-angle et de la vitesse constitutive fournit des expressions
équivalentes de la même évaluation \cite{hmtg26metric,hmtg26counter}.
La section longitudinale de référence, l’action et la vitesse
dimensionnelle restent explicites. Le choix \(L_*=1\,\mathrm m\)
est la section utilisée pour les chiffres SI ; sa déclaration permet
de les reproduire avec le même sens dimensionnel. L’identification
physique R8 conserve sa place dans la démonstration.
